Recomendaciones generales para el formato ICPC. Conviene confirmar los detalles (duración, penalización, congelamiento del marcador) en las reglas del concurso concreto.

\begin{props}
  \item \textbf{El computador es el recurso escaso.} Normalmente el equipo de 3 comparte un solo computador. Quien no está programando piensa o escribe a mano; se pasa al teclado cuando el algoritmo ya está decidido.
  \item \textbf{Primeros minutos:} repartir los problemas para que entre los tres se lean todos, y clasificarlos por dificultad. El marcador ayuda: los problemas que resuelven más equipos suelen ser los más fáciles.
  \item \textbf{Empezar por lo seguro.} Resolver primero los problemas fáciles, rápido pero sin descuidar la verificación.
  \item \textbf{Cada envío rechazado cuesta.} El ICPC suele sumar una penalización de tiempo (normalmente 20 minutos) por cada envío rechazado de un problema que luego se resuelve. Antes de enviar: ejemplos del enunciado, casos borde (\(n = 1\), vacío, todos iguales, valores máximos) y un caso del tamaño máximo para detectar TLE.
  \item \textbf{Ante un Wrong Answer:} releer el enunciado, revisar los límites y tipos, y pedir que otro integrante lea el código (idealmente impreso, mientras el computador lo usa alguien más). Si no se ve el error, usar la prueba de estrés contra una fuerza bruta.
  \item \textbf{No quedarse atascado.} Si un problema no avanza tras un tiempo razonable (por ejemplo, 30 minutos), rotarlo a otro integrante o cambiar de problema.
  \item \textbf{Usar el notebook con criterio:} copiar el algoritmo no sustituye entender los tipos de datos, los índices (0 o 1) y los límites del problema.
  \item \textbf{Leer las aclaraciones publicadas} por el jurado: pueden cambiar la interpretación de un enunciado.
  \item \textbf{Última hora:} el marcador suele congelarse; conviene dejar de lado los problemas casi imposibles y asegurar envíos correctos de los que sí se pueden terminar.
\end{props}
